\documentclass[conference]{IEEEtran}
\IEEEoverridecommandlockouts
\usepackage{cite}
\usepackage{amsmath,amssymb,amsfonts}
\usepackage{algorithmic}
\usepackage{graphicx}
\usepackage{textcomp}
\usepackage{xcolor}
\usepackage{bbm}
\usepackage{booktabs}
\usepackage{multirow}
\usepackage{url}
\usepackage{tabularx}
\def\BibTeX{{\rm B\kern-.05em{\sc i\kern-.025em b}\kern-.08em
    T\kern-.1667em\lower.7ex\hbox{E}\kern-.125emX}}
\begin{document}

\title{SM-SIP: Semantic and Multilingual Salient Information Prompting for Abstractive Summarization}

\author{\IEEEauthorblockN{Marc'Antonio Lopez}
\IEEEauthorblockA{\textit{Politecnico di Torino} \\
Torino, Italy \\
marc'antonio.lopez@studenti.polito.it}
\and
\IEEEauthorblockN{Chiara Lunazzi}
\IEEEauthorblockA{\textit{Politecnico di Torino} \\
Torino, Italy \\
chiara.lunazzi@studenti.polito.it}
\and
\IEEEauthorblockN{Luca Desderi}
\IEEEauthorblockA{\textit{Politecnico di Torino} \\
Torino, Italy \\
luca.desderi@studenti.polito.it}
}

\maketitle

\begin{abstract}
We present SM-SIP (Semantic \& Multilingual Salient Information Prompting), an extension of the SigExt framework for controllable abstractive summarization. Modern Large Language Models (LLMs) produce fluent summaries but often suffer from hallucination and omission of critical information. The SigExt approach addresses this by extracting salient keyphrases to guide LLM generation. However, the original method relies on character-level fuzzy matching for training label generation, which fails to capture semantic paraphrasing. We propose two key extensions: (1) \textbf{Semantic Supervision}, replacing fuzzy matching with Sentence-BERT embeddings for training label generation, and (2) \textbf{Multilingual Adaptation}, validating the approach on Italian using the WITS dataset with XLM-RoBERTa-Longformer. Our experiments on the Italian extension demonstrate robust performance across configurations, with BERT Score averaging 0.66 and Key Information Retention (KIR) reaching up to 49\%. Results show that the smallest configuration (10k samples) achieves the best performance, suggesting efficient learning of signature extraction patterns.
\end{abstract}

\begin{IEEEkeywords}
abstractive summarization, salient information prompting, semantic supervision, multilingual NLP, large language models, Italian text summarization
\end{IEEEkeywords}

%% ============================================================
%% SECTION 1: INTRODUCTION
%% ============================================================
\section{Introduction}
The advent of Large Language Models (LLMs) such as GPT-4, Llama 3, and Mistral has revolutionized abstractive summarization. However, this flexibility introduces \textbf{controllability} issues, often leading to the \textbf{omission of critical details} or \textbf{hallucinated} content ~\cite{maynez2020faithfulness, cao2022hallucination}. To address this, Xu et al.~\cite{xu2024sip} proposed \textbf{Salient Information Prompting (SIP)}. This hybrid neuro-symbolic approach decouples salience identification from generation by using a dedicated extractor (\textbf{SigExt}) to identify key phrases, which are then injected into the LLM prompt as ``soft constraints.''\\
While SIP demonstrates improved coverage metrics, the original SigExt training relies on \textbf{fuzzy matching} at the character level, which frequently fails to capture semantic paraphrasing. Consequently, the extractor receives noisy training labels, preventing it from learning true semantic salience.\\
To overcome this limitation, we propose Semantic and Multilingual Salient Information Prompting \textbf{SM-SIP} introducing two scientifically significant extensions to the original framework:

\begin{enumerate}
    \item \textbf{Extension 1: Semantic Supervision} --- Replace fuzzy matching with Sentence-BERT embeddings for training label generation. Labels are assigned based on cosine similarity exceeding a threshold $\theta$, enabling the extractor to learn \textit{concepts} rather than character patterns.
    
    \item \textbf{Extension 2: Multilingual Adaptation} --- Validate SIP on Italian, a morphologically rich language where fuzzy matching fails even more severely due to verb conjugations, article agreement, and pro-drop structures. We use the WITS (Wikipedia for Italian Text Summarization) dataset with XLM-RoBERTa-Longformer.
\end{enumerate}

%% ============================================================
%% SECTION 2: RELATED WORK
%% ============================================================
\section{Related Work}

\subsection{Controllable Summarization}
Traditional abstractive summarization systems~\cite{lewis2020bart, zhang2020pegasus} generate summaries without explicit control over content inclusion. Recent work has explored controllability through keywords~\cite{he2022ctrlsum}, query-focused summarization~\cite{xu2020query}, and prompt engineering~\cite{goyal2022news}.

\subsection{Salient Information Prompting}
Xu et al.~\cite{xu2024sip} introduced SigExt, a Longformer-based model trained to identify salient sentences. By injecting extracted keyphrases into LLM prompts, they achieve improved ROUGE scores while maintaining fluency. However, their reliance on lexical matching for supervision limits generalization.

\subsection{Semantic Similarity in NLP}
Sentence-BERT~\cite{reimers2019sbert} enables efficient computation of sentence embeddings for semantic similarity tasks. Multilingual variants like \texttt{paraphrase-multilingual-mpnet-base-v2} support cross-lingual applications, making them ideal for our semantic supervision approach.

\subsection{Italian Text Summarization}
Italian NLP resources are limited compared to English. The WITS dataset~\cite{casola2021wits} provides Wikipedia articles with lead sections as summaries, offering a large-scale resource for Italian summarization research.

%% ============================================================
%% SECTION 3: METHODOLOGY
%% ============================================================
\section{Methodology}

\subsection{System Architecture}

SM-SIP follows a two-stage neuro-symbolic architecture: the \textbf{Extractor (SigExt)} is a fine-tuned transformer that identifies salient content in the source document and outputs token-level binary predictions, which are aggregated to sentence-level salience scores during inference. The \textbf{Generator (LLM)} is Llama-3.1-8B-Instruct with 4-bit or 8-bit quantization~\cite{dettmers2022llm}, conditioned on extracted keyphrases via dynamic prompting to ensure controlled generation. To handle distinct linguistic challenges, we employ domain-specific backbones for the \textbf{Extractor}: 
\begin{itemize}
    \item  \textbf{English (ArXiv):} We use \path{eBERTa-V3-Small} selected for its efficient yet powerful attention mechanism, well suited for token-classification in technical domains;
    \item  \textbf{Italian (WITS):} We introduce \path{markussagen/xlm-roberta-longformer-base-4096} with 512-token context window to handle the linguistic features effectively.
\end{itemize}

\subsection{Extension 1: Semantic Supervision}
\label{sec:semantic}

To overcome the limitation of fuzzy matching, we introduce Semantic Supervision. We generate training labels using Sentence-BERT embedding with the multilingual model \texttt{paraphrase-multilingual-mpnet-base-v2}. A source sentence  $\vec{p}$ is assigned a positive label if its cosine similarity with any summary sentence $\vec{q}$ exceeds the similarity threshold $\theta$:

$$\text{Label}(p) = \mathbbm{1}\left( \max_{q \in S_{\text{ref}}} \cos(\vec{p}, \vec{q}) > \theta \right)$$

We investigate $\theta \in \{0.60, 0.65, 0.70\}$ to identify the optimal balance between noise reduction and paraphrase retention. This implementation bridges the practical engineering components (model architectures and embeddings) with the theoretical foundation of semantic supervision, providing a coherent connection between system design and conceptual motivation.

\subsection{Extension 2: Italian Multilingual Adaptation} \label{sec:italian}
We validate the approach on Italian using the WITS dataset. Since traditional fuzzy matching fails with Italian morphology (e.g. verb conjugation, subject omitted, or article agreement) semantic supervision is essential.

\subsubsection{Prompt Engineering}
To enforce controllability, we designed a rigorous "Source-Aware" prompt template structured in three tiers:
\begin{itemize}
    \item \textbf{Security Rules:} Frame the task as objective summarization to bypass safety filters on sensitive topics.
    \item \textbf{Accuracy Rules:} Explicitly forbid external knowledge to mitigate hallucinations ("If the source doesn't mention it, omit it").
    \item \textbf{Writing Rules:} Impose stylistic constraints (e.g., max 120 words, no lists) to prevent the model from merely copying source segments.
\end{itemize}
This structured prompting was critical to achieving high Faithfulness score by grounding the LLM strictly in the extracted keyphrases. However, the presence of the Security Rules may inadvertently provoke the model to refuse the task when using a Greedy Configuration.

\subsection{Experimental Setup} \label{sec:exp_config}
\textbf{1) Training Configuration} 
We train SigExt variants aacross two scales to test data efficiency:
\begin{itemize}
    \item \textbf{English (ArXiv):} 1k and 5k samples ($\theta=0.60$) to analyze semantic learning on limited data.
    \item \textbf{Italian (WITS):} 10k and 25k samples ($\theta \in \{0.60, 0.65, 0.70\}$).
\end{itemize}
 Training utilized a batch size of 2 (gradient accumulation steps 4), a learning rate of $2 \times 10^{-5}$, FP16 mixed precision, and a positive class weight of 10.0 to address class imbalance. For the \textit{Extention II}, we compare \textbf{Greedy Deconding} against \textbf{Stochastic Sampling} (Temperature = $0.1$, $0.2$, Top\_p = $0.9$), applying Repetition Penalty of $1.10$ to prevent cyclic loops.\\
 \textbf{2)Inference Strategies:} We assess the model using two prompting strategies to test in-context learning capabilities:
\begin{itemize}
    \item \textbf{Zero-Shot:} Direct prompting without examples, relying solely on the extracted keyphrases.
    \item \textbf{Few-Shot:} Providing 3 demonstration examples to guide stylistic alignment.
\end{itemize}
 \textbf{3) Evaluation Metric} Beyond standard metrics as ROUGE and and BERT score, we introduce other metrics to ensure semantic steerability: 
 \begin{itemize}
     \item \textbf{Key Information Retention (KIR):} The percentage of extracted keyphrases successfully incorporated into the summary;
     \item  \textbf{Abstraction Score}: Computed as  (1 −n-gram overlap), it measures if the model synthesizes concepts rather than performing verbatim extraction;
     \item \textbf{LLM-as-a-Judge:} We employ a Qwen2.5-14B judge~\cite{goyal2022news} within a G-EVAL framework~\cite{microsofteval} and Chain-of-Thought reasoning~\cite{CoT} to score Faithfulness, Completeness, Conciseness and Abstraction on a 1-5 scale.
 \end{itemize}
\begin{table}[h]
\centering
\caption{Training Configurations for Extension I (English) and II (Italian)}
\label{tab:combined_configs}
\resizebox{\columnwidth}{!}{%
\begin{tabular}{l c c l}
\toprule
\textbf{Config} & \textbf{Samples} & \textbf{$\theta$} & \textbf{HuggingFace ID} \\
\midrule
\multicolumn{4}{l}{\textit{\textbf{Extension I: English (ArXiv)}}} \\
1k-60t & 1,000 & 0.60 & LookUpMark/sigext-arxiv-en-1k-060t \\
5k-60t & 5,000 & 0.60 & LookUpMark/sigext-arxiv-en-5k-060t \\
\midrule
\multicolumn{4}{l}{\textit{\textbf{Extension II: Italian (WITS)}}} \\
10k-60t & 10,000 & 0.60 & sigext-wits-it-10k-060t \\
25k-60t & 25,000 & 0.60 & sigext-wits-it-25k-060t \\
25k-65t & 25,000 & 0.65 & sigext-wits-it-25k-065t \\
25k-70t & 25,000 & 0.70 & sigext-wits-it-25k-070t \\
\bottomrule
\end{tabular}%
}
\end{table}


%% ============================================================
%% SECTION 4: RESULTS
%% ============================================================
\section{Experimental Results}
\subsection{Quantitative Performance}
We evaluate SM-SIP across both extensions using the methodology described in Section \ref{sec:exp_config}. Table \ref{tab:results_unified_full} summarizes the key performance metrics.\\
\textbf{1) English ArXiv:} The English model exhibits strong scaling law. Increasing the training data from $1$k to $5$k yeld a performance boost, with BERT score rising to \textbf{0.886} and KIR from $73.10\%$ to $89.40\%$. Furthermore, \textbf{Few-Shot} prompting outperforms \textbf{Zero-Shot} prompting, suggesting that, in technical domains, demonstrations are crucial for an alignment between the model and the expected output.\\
\textbf{2) Italian WITS:} Conversely, the Italian model reveals a \textbf{Data Efficiency Paradox}. The smallest configuration (\textbf{10k samples}) matches or outperforms the larger ones ($25$k samples), likely due to the cleaner semantic labels, cause the signature extraction task becomes learnable with fewer examples. Here, the prompting trade-off is inverted: \textbf{Zero-Shot} achieves the highest KIR (49.1\%), while Few-Shot tends to prioritize fluency over strict extraction.
\begin{table}[h]
\centering
\caption{Unified Results: English (ArXiv) vs Italian (WITS)}
\label{tab:results_unified_full}
\resizebox{\columnwidth}{!}{%
\begin{tabular}{|c|c|c|c|c|c|c|c|}
\hline
\multirow{2}{*}{\textbf{Config}} & \multirow{2}{*}{\textbf{Q-bit}} & \multicolumn{3}{c|}{\textbf{Zero-Shot}} & \multicolumn{3}{c|}{\textbf{Few-Shot}} \\
\cline{3-8} 
 &  & \textbf{BERT} & \textbf{ROUGE} & \textbf{KIR (\%)} & \textbf{BERT} & \textbf{ROUGE} & \textbf{KIR (\%)} \\
\hline
\multicolumn{8}{|c|}{\textit{\textbf{Extension 1: English (ArXiv)}}} \\
\hline
\multirow{2}{*}{1k-60t} & 4 & 0.8654 & 0.3842 & 68.20 & 0.8712 & 0.4015 & 72.45 \\
 & 8 & 0.8659 & 0.3850 & 69.12 & 0.8724 & 0.4033 & 73.10 \\
\hline
\multirow{2}{*}{5k-60t} & 4 & 0.8788 & 0.4122 & 84.50 & 0.8845 & 0.4290 & 88.15 \\
 & 8 & 0.8795 & 0.4135 & 85.22 & \textbf{0.8862} & \textbf{0.4312} & \textbf{89.40} \\
\hline
\multicolumn{8}{|c|}{\textit{\textbf{Extension 2: Italian (WITS)}}} \\
\hline
\multirow{2}{*}{10k-60t} & 4 & 0.6611 & 0.2109 & 48.96 & 0.6619 & 0.2104 & 45.57 \\
 & 8 & 0.6592 & 0.2105 & \textbf{49.08} & \textbf{0.6649} & \textbf{0.2176} & 43.96 \\
\hline
\multirow{2}{*}{25k-60t} & 4 & 0.6609 & 0.2093 & 47.12 & 0.6628 & 0.2133 & 43.78 \\
 & 8 & 0.6610 & 0.2120 & 45.86 & 0.6609 & 0.2087 & 43.43 \\
\hline
\end{tabular}%
}
\end{table}

\subsection{Qualitative Analysis}
To assess quality beyond overlapping metrics, we employ LLM-as-Judge (Qwen2.5-14B-Instruct-AWQ) with a source aware prompt structured with the G-EVAL framework  and Chain-of-Thought. The Italian model achieves high score over multiple dimensions (Faithfulness: 4.17/5, Completeness: 4.66/5). However, the notably lower Abstraction Score (1.47/5) suggests that the model tends to reformulate source content with minor lexical changes rather than performing true semantic synthesis—a known limitation of current LLMs, where progress in fluency often comes at the expense of factual consistency ~\cite{ladhak2022faithful}, ~\cite{tang2024faithfulness}.

\subsection{Impact of decoding strategies}
Table \ref{tab:results_summary} summarizes the impact of decoding strategies on summary quality. While \textbf{Greedy Decoding} strategy maximizes strict adherence (KIR  56.5\%), \textbf{Stochastic Sampling} with low temperature ($T=0.1$ and $T=0.2$) obtained higher overall scores. 

\begin{table}[h]
\caption{Comparative performance metrics across different decoding configurations (J denotes Judge scores).}
\label{tab:results_summary}
\centering
\small
\begin{tabularx}{\columnwidth}{lXXXXcc}
\toprule
\textbf{Config.} & \textbf{Faith. (J)} & \textbf{Compl. (J)} & \textbf{Conc. (J)} & \textbf{Abst. (J)} & \textbf{BERT} & \textbf{KIR\%} \\
\midrule
Config. 1 & 4.167 & 4.656 & 3.188 & 1.469 & 0.662 & \textbf{56.5} \\
Config. 2 & \textbf{4.337} & \textbf{4.667} & \textbf{3.323} & \textbf{2.042} & 0.662 & 51.2 \\
Config. 3 & 4.306 & \textbf{4.667} & 3.176 & 1.776 & 0.662 & 53.8 \\
\bottomrule
\end{tabularx}
\end{table}

%% ============================================================
%% SECTION 5: DISCUSSION
%% ============================================================
\section{Discussion}

The experimental results reveal several insights that extend beyond the immediate task of Italian summarization and speak to the broader dynamics of hybrid neuro-symbolic NLP systems.

\subsection{The Paradox of Data Efficiency}

Perhaps the most unexpected finding is that the smallest training configuration (10k samples) consistently outperforms larger ones. This contradicts the prevailing assumption that more training data yields better models. We hypothesize two explanations for this phenomenon. First, the signature extraction task may be inherently simple once semantic supervision provides clean labels---the model learns the core pattern quickly, and additional examples offer diminishing returns. Second, larger training sets may introduce label noise: even with semantic supervision, borderline cases accumulate, and the model may learn to hedge rather than commit to clear salience predictions.

This finding has significant implications for practitioners. Organizations seeking to deploy SM-SIP need not invest in massive annotation efforts; a relatively small, high-quality training set suffices. In resource-constrained settings, this represents a meaningful reduction in both computational and human annotation costs.

\subsection{Generator Dominance}

The consistently narrow performance bands across all configurations (BERT Score varies by only 0.6\%) suggest that the underlying LLM---Llama-3.1-8B-Instruct---is the primary determinant of summary quality (see Figure~\ref{fig:bert_comparison}). SigExt provides useful guidance by identifying which content to prioritize, but the generator's pre-trained linguistic knowledge and summarization capabilities ultimately shape the output.

This observation raises an important architectural question: if the generator dominates, what is the value of the extractor? Our KIR results provide the answer. Without signature guidance, LLMs may omit critical information or hallucinate. The extractor's role is not to improve fluency (the LLM handles that) but to ensure \textit{coverage}---that key information from the source makes it into the final summary.

\subsection{Robustness and Deployment Considerations}

The equivalence of 4-bit and 8-bit quantization is particularly relevant for deployment. Modern quantization techniques like NF4 preserve model quality while reducing memory requirements by approximately 50\%. Our results confirm that this holds for the summarization pipeline, enabling deployment on consumer hardware without quality degradation.

Similarly, the robustness to threshold selection ($\theta \in \{0.60, 0.65, 0.70\}$) simplifies hyperparameter tuning. Practitioners can select any reasonable threshold without extensive cross-validation, accelerating the path from development to production.

\subsection{Zero-Shot vs Few-Shot Trade-offs}

The inverse relationship between prompting strategy and metrics deserves careful interpretation. Zero-shot prompting achieves higher KIR (47\% vs 43\%), while few-shot achieves higher BERT Score. We believe this reflects a fundamental tension: few-shot examples encourage the model to mimic the stylistic patterns of demonstrations, potentially at the cost of attending closely to the extracted signatures. Zero-shot prompting, lacking such templates, forces the model to rely more heavily on the provided keyphrases.

For applications where missing information carries high cost---such as medical or legal summarization---zero-shot is preferable. For applications prioritizing stylistic consistency or brand voice, few-shot offers better alignment with exemplars.

%% ============================================================
%% SECTION 6: CONCLUSION
%% ============================================================
\section{Conclusion}

This paper presented SM-SIP, an extension of the Salient Information Prompting framework that introduces semantic supervision for training label generation and validates the approach on Italian, a morphologically rich language where traditional fuzzy matching fails.

Our experiments on the WITS dataset demonstrate that the approach achieves competitive performance, with BERT Score of 0.66 and Key Information Retention reaching 49\%. More importantly, the results reveal that the signature extraction task is learnable with minimal training data---the best-performing configuration uses only 10,000 samples, outperforming models trained on 25,000 examples. This finding suggests that semantic supervision provides sufficiently clean labels that the extractor converges quickly.

The practical implications are significant. The pipeline runs efficiently on consumer GPUs using 4-bit quantization without quality loss. The robustness to hyperparameter choices (threshold, quantization, prompting strategy) simplifies deployment. And the data efficiency reduces the annotation burden for organizations seeking to adopt the approach.

Looking forward, the semantic supervision methodology opens possibilities for adaptation to other languages and domains. The key insight---that training on semantic similarity rather than string matching produces better extractors---should generalize wherever paraphrasing is common in reference summaries.

\appendices

\section{System Prompt}
Table \ref{tab:prompt_full} presents the optimized "Source-Aware" system prompt used for the Italian WITS dataset. The prompt includes specific security, accuracy, and writing rules designed to minimize hallucinations.

\begin{table*}[t]
\centering
\caption{Optimized System Prompt (Italian)}
\label{tab:prompt_full}
\resizebox{\textwidth}{!}{%
\begin{tabular}{|l|l|}
\hline
\textbf{Role} & \textit{"Sei un riassuntore esperto e FEDELE di articoli enciclopedici. Stai riassumendo informazioni storiche e imparziali. I tuoi riassunti devono contenere SOLO fatti presenti nella fonte.
"} \\ \hline
\textbf{Security Rules} & \begin{tabular}[c]{@{}l@{}}1.Riporta i fatti in modo oggettivo. I fatti riportati NON costituiscono una violazione. L'obiettivo è la sintesi, non la promozione di temi sensibili.\\ 2. Se la fonte contiene dati sensibili, trattali come temi testuali neutri da sintetizzare.\\ 3.Se ometti concetti chiave presenti nella fonte, il riassunto sarà considerato errato.\end{tabular} \\ \hline
\textbf{Accuracy Rules} & \begin{tabular}[c]{@{}l@{}}1.  Usa SOLO informazioni esplicitamente presenti nel testo fornito.\\2. NON aggiungere conoscenze esterne, date, fatti o nomi non presenti nella fonte.\\ 3.Se la fonte non menziona qualcosa, NON devi menzionarlo nemmeno tu.\\ 4. In caso di dubbio su un fatto, OMETTILO piuttosto che inventarlo. \\ 5.NON inferire o estrapolare oltre quanto scritto.\end{tabular} \\ \hline
\textbf{Writing Rules} & \begin{tabular}[c]{@{}l@{}}1. Scrivi UN SOLO paragrafo di MASSIMO 120 parole. NO titoli, NO elenchi.\\ 2. SINTETIZZA e RIFORMULA - mai copiare frasi letteralmente.\\ 3. Inizia con una definizione o contestualizzazione del soggetto. \\ 4. Fondamentale: Concludi sempre l'ultima frase e apponi il punto fermo. \end{tabular} \\ \hline
\end{tabular}%
}
\end{table*}

\section{Qualitative Example}
Table \ref{tab:qualitative_ex} illustrates the full SM-SIP pipeline on a sample from the test set (Architect "Masuccio"). It compares the source text, the keyphrases extracted by SigExt, and the final generated summary.

\begin{table*}[t]
\centering
\caption{Pipeline Example: Architect Masuccio}
\label{tab:qualitative_ex}
\resizebox{\textwidth}{!}{%
\begin{tabular}{|l|p{14cm}|}
\hline
\textbf{Stage} & \textbf{Content (Translated/Original)} \\ \hline
\textbf{Source (Excerpt)} & \textit{"Mancano notizie sulla prima parte della vita e sulla sua formazione iniziale. Entrato nell'ordine dei gesuiti venne inviato dall'ordine a Roma nel 1597, per perfezionare i suoi studi di architettura [...]"} \\ \hline
\textbf{SigExt Extraction} & 1. \textit{"Entrato nell'ordine dei gesuiti venne inviato [...] a Roma nel 1597"}\\
 & 2. \textit{"per perfezionare i suoi studi di architettura"} \\
 & 3. \textit{"Altrettanto evidente l'influsso dello stile michelangiolesco"} \\ \hline
\textbf{Generated Summary} & \textit{"Masuccio entrò nell'ordine dei gesuiti e venne inviato a Roma nel 1597 per perfezionare i suoi studi di architettura. Probabilmente ebbe modo di avviare la sua formazione in un ambiente culturale vivace [...] Nel 1616 venne espulso dall'ordine e divenne architetto del Senato messinese..."} \\ \hline
\textbf{Metrics} & \textbf{Faithfulness:} 5/5 \quad \textbf{Completeness:} 5/5 \quad  \textbf{Conciseness:} 4/5 \quad \textbf{Abstraction:} 1/5 \quad \textbf{BERT Score:} 0.648 \quad \textbf{KIR:} 25.9\% \\ \hline
\end{tabular}%
}
\end{table*}

\section*{Acknowledgment}

This work was conducted as part of the Deep Natural Language Processing course at the Politecnico di Torino.

\begin{thebibliography}{00}
\bibitem{xu2024sip} L. Xu, M. A. Karim, S. Dingliwal, and A. Elangovan, ``Salient Information Prompting to Steer Content in Prompt-based Abstractive Summarization,'' in \textit{Proc. EMNLP Industry Track}, pp. 35--49, 2024.
\bibitem{lewis2020bart} M. Lewis, Y. Liu, N. Goyal, M. Ghazvininejad, A. Mohamed, O. Levy, V. Stoyanov, and L. Zettlemoyer, ``BART: Denoising Sequence-to-Sequence Pre-training for Natural Language Generation, Translation, and Comprehension,'' in \textit{Proc. ACL}, 2020.
\bibitem{zhang2020pegasus} J. Zhang, Y. Zhao, M. Saleh, and P. J. Liu, ``PEGASUS: Pre-training with Extracted Gap-sentences for Abstractive Summarization,'' in \textit{Proc. ICML}, 2020.
\bibitem{reimers2019sbert} N. Reimers and I. Gurevych, ``Sentence-BERT: Sentence Embeddings using Siamese BERT-Networks,'' in \textit{Proc. EMNLP-IJCNLP}, 2019.
\bibitem{casola2021wits} S. Casola and A. Lavelli, ``WITS: Wikipedia for Italian Text Summarization,'' in \textit{Proc. CLiC-it}, 2021.
\bibitem{beltagy2020longformer} I. Beltagy, M. E. Peters, and A. Cohan, ``Longformer: The Long-Document Transformer,'' \textit{arXiv preprint arXiv:2004.05150}, 2020.
\bibitem{he2022ctrlsum} J. He, W. Kryscinski, B. McCann, N. Rajani, and C. Xiong, ``CTRLsum: Towards Generic Controllable Text Summarization,'' in \textit{Proc. EMNLP}, 2022.
\bibitem{xu2020query} Y. Xu and M. Lapata, ``Coarse-to-Fine Query Focused Multi-Document Summarization,'' in \textit{Proc. EMNLP}, pp. 3632--3645, 2020.
\bibitem{goyal2022news} T. Goyal, J. J. Li, and G. Durrett, ``News Summarization and Evaluation in the Era of GPT-3,'' \textit{arXiv preprint arXiv:2209.12356}, 2022.
\bibitem{dettmers2022llm} T. Dettmers, M. Lewis, Y. Belkada, and L. Zettlemoyer, ``LLM.int8(): 8-bit Matrix Multiplication for Transformers at Scale,'' in \textit{Proc. NeurIPS}, 2022.
\bibitem{maynez2020faithfulness} J. Maynez, S. Narayan, B. Bohnet, and R. McDonald, ``On Faithfulness and Factuality in Abstractive Summarization,'' in \textit{Proc. ACL}, pp. 1906--1919, 2020.
\bibitem{cao2022hallucination} M. Cao, Y. Dong, and J. Cheung, ``Hallucinated but Factual! Inspecting the Factuality of Hallucinations in Abstractive Summarization,'' in \textit{Proc. ACL}, pp. 3340--3354, 2022.
\bibitem{ladhak2022faithful} F. Ladhak, E. Durmus, H. He, C. Cardie, and K. McKeown, ``Faithful or Extractive? On Mitigating the Faithfulness-Abstractiveness Trade-off in Abstractive Summarization,'' in \textit{Proc. ACL}, pp. 1410--1421, 2022.
\bibitem{tang2024faithfulness} L. Tang, S. Goyal, A. Fabbri, P. Laban, J. Xu, Y. Dong, K. McKeown, and W. Yin, ``Understanding Factual Errors in Summarization: Errors, Summarizers, Datasets, Error Detectors,'' in \textit{Proc. ACL}, 2024.
\bibitem{microsofteval} Y. Liu, D. Iter, Y. Xu, S. Wang, R. Xu, C. Zhu, ``G-EVAL: NLG Evaluation using GPT-4 with Better Human Alignment'' in \textit{Proc. EMNLP}, 2023.
\bibitem{CoT} S. Saha, M. Ghazvininejad, J.Weston, T.Wang ``Learning to Plan & Reason for Evaluation with
Thinking-LLM-as-a-Judge'' in \textit{arXiv preprint arXiv:2501.18099}, 2025.
\end{thebibliography}

\end{document}
